\documentclass[10pt,a4paper]{amsart}
\usepackage[margin=19mm]{geometry}
\usepackage{amsmath,amssymb,mathtools}
\usepackage[scr=rsfs,cal=euler]{mathalfa}
\usepackage{xfrac}
\usepackage[unicode,hidelinks,bookmarks=false]{hyperref}
\newcommand{\ie}{i.e.\ }
\newcommand{\cf}{cf.\ }
\newcommand{\resp}{respectively\ }
\newcommand{\Subsection}{\subsection}
\providecommand{\nobreakdash}{\nobreak\hbox{-}}
\setlength{\emergencystretch}{2em}
\multlinegap0pt
\usepackage{float}
\usepackage{amssymb}
\usepackage{xcolor}
\usepackage[shortlabels]{enumitem}
\newtheorem{cor}{Corollary}
\newtheorem{lem}{Lemma}
\newtheorem{thm}{Theorem}
\newtheorem{prop}{Proposition}
\title{Newton polygons for the non-bipartite dimer model}
\author{Vladimir Bo\v{s}kovi\'c}
\date{Source: arXiv:2607.07503v1 (CC BY 4.0)}
\begin{document}
\maketitle

\noindent\emph{Source and licence.} Newton polygons for the non-bipartite dimer model. Version arXiv:2607.07503v1 (CC BY 4.0), distributed by arXiv under the Creative Commons Attribution 4.0 International licence, http://creativecommons.org/licenses/by/4.0/, as recorded in the arXiv licence metadata for this exact version and shown on the abstract page https://arxiv.org/abs/2607.07503v1. Full licence notice: \texttt{licenses/arxiv-2607.07503-CC-BY-4.0.txt}.\par\smallskip

\noindent\textit{Editorial scope.} Counted unit: the article's thm marginals (source0.tex lines 814--816). The article's complete original proof of it is delivered with it (source0.tex lines 818--847, one \texttt{proof} environment of 30 lines). No mathematics is written, repaired or re-derived: the statement and the proof are the article's own lines, in the article's own notation, and no retained line is substituted or re-flowed.  Retained uncounted as the source-local context the proof needs: lines 320--322; lines 379--381; lines 396--400; lines 481--484; lines 502--506; lines 691--695; lines 740--747.  One declared adaptation: exactly 1 ancillary strings inside the retained spans are omitted (image-inclusion commands and, where the source repeats a label, the repeated label definitions) and no other line differs from the source; the figure environments, with their placement, captions and labels, are kept, and the package records the omitted strings in fidelity receipts bound to the affected bodies.  No retained line contains a citation, so the wrapper carries no bibliography and no bibliography entry.  The retained lines contain the article's own diagram code, which the wrapper typesets with the standard tikz packages; no image file is delivered and the package declares no asset dependency.  Editorial formatting only, in addition: an AMS \texttt{amsart} preamble that restates the article's own declarations and macros used by the retained lines, namely the environments \texttt{cor}, \texttt{lem}, \texttt{thm}, \texttt{prop} on separate counters, together with the standard packages those lines need.  The retained environments print in this document as Corollary 1, Lemma 1, Theorem 1, Proposition 1 in order of appearance, whereas the article prints the same items with the numbering of its own full document.  The complete native source of the article as distributed in the arXiv e-print for this exact version is bundled unchanged as \texttt{sources/204713/source0.tex}.
\begin{cor} \label{NPs and homology classes}
    Let $\Gamma$ be a minimal bipartite torus graph, and let $P(z,w)$ be its characteristic polynomial obtained from the full Kasteleyn matrix. There is a bijection between the multiset of homology classes of the zig-zag paths on $\Gamma$ and the primitive edge vectors on the Newton polygon $N(P)$. 
\end{cor}

\begin{lem} \label{vertical zig-zag path}
    Let $\Gamma$ be a torus graph that contains a zig-zag path $Z$ with a homology $(0,1)$. Then we can choose a fundamental domain, without changing the homology classes of the zig-zag paths of $\Gamma$, such that the vertical boundary of the fundamental domain crosses every edge of $Z$ once and no other edges of $\Gamma$. 
\end{lem}

\begin{lem} \label{fundamental_domain matching}
    Let $\Gamma$ be an isoradial graph.  
    Suppose the homology classes of its zig-zag paths are given by vectors $[Z_i] = (\alpha_i, \beta_i)$ for $1 \leq i \leq n$, and assume $[Z_1] = (0,1)$. Define \[k = \sum_{i=1}^n \max(\alpha_i,0).\] 
    Then the zig-zag path $Z_1$ consists of exactly $k$ edges.
\end{lem}

\begin{lem} \label{N(Gamma) < N(Gamma_Delta_v)}
    Let $\Gamma$ be a torus graph containing a $3$-valent vertex $v \in V(\Gamma)$, and let $\Gamma_{\Delta_v}$ be the graph obtained by replacing the vertex $v$ with a triangle. Then
    \[N(\Gamma)\subseteq N(\Gamma_{\Delta_v}).\]
\end{lem}

\begin{thm} \label{equal newton polygons delta}
    Let $\Gamma$ be an isoradial bipartite graph. Then
    \[N(\Gamma_{\Delta_{S}}) = N(\Gamma),\]
    where $\Gamma_{\Delta_{S}}$ denotes the graph obtained by replacing the vertices $S \subseteq V_3 = \{v \in V(\Gamma): \deg(v) = 3\}$ with triangles.
\end{thm}

\begin{prop} \label{rootsT}
    Let $m,n \in \mathbb{N}$ and assume that $m$ is even. The polynomials $P_R(w)$ and $P_U(z)$ have all real roots, and each root can be written as 
    \[r_i = \prod_{1 \leq j \leq m} \frac{a_{j,i}}{b_{j,i}}, \text{ for } 1 \leq i \leq n,\]
    \[u_j = (-1)^{(j-1)n} \prod_{1 \leq i \leq n} \frac{c_{j,i}}{b_{j,i}} ,\text{ for } 1 \leq j \leq m.\]
\end{prop}

\begin{lem} \label{no double edges}
    Let $\gamma \in \mathcal{L}(F_{m,n})$ be a collection of loops with a realizable homology $(m,i)$ for any $-n \leq i \leq 0$. Then $\gamma$ satisfies the following properties:
    \begin{enumerate} [(a)]
        \item The non-trivial loops in $\gamma$ traverse all $b$-edges of $F_{m,n}$;
        \item The only trivial loops in $\gamma$ are double edges traversing an $a$- or $c$-edge;
        \item For every column of $F_{m,n}$, the non-trivial loops in $\gamma$ traverse either all $a$-edges or all $c$-edges.
    \end{enumerate}
\end{lem}

\begin{thm} \label{real rooted marginals}
    Let $\Gamma$ be an isoradial bipartite torus graph, and let $S \subseteq V(\Gamma)$ be a subset of vertices of degree $3$. Then the marginal polynomials of $N(\Gamma_{\Delta_S})$ are real-rooted.
\end{thm}

\begin{proof}
    Without loss of generality, up to an $SL_2(\mathbb{Z})$, we can assume that $N(\Gamma)$ contains a vertical side. Since $\Gamma$ is isoradial bipartite, we can apply Corollary \ref{NPs and homology classes} to conclude that $N(\Gamma)$ is described by the homology classes of the zig-zag paths in $\Gamma$. Therefore, we can assume that $\Gamma$ has $r \geq 1$ zig-zag paths $\{Z_1, \dots, Z_r\}$ with homology class $(0,1)$.

    Using Lemma \ref{vertical zig-zag path}, the graph $\Gamma$ can be embedded in the torus without changing its characteristic polynomial, such that there is a zig-zag path, say $Z_r$, where the vertical boundary, denoted by $\gamma_y$, of the fundamental domain crosses all the edges of $Z_r$ and no other edges of $\Gamma$. As discussed in Lemma \ref{fundamental_domain matching}, each zig-zag path $Z_i$, $1 \leq i \leq r$ has exactly $k$ vertices, where $k$ describes the width of $N(\Gamma)$. We also fix the horizontal boundary $\gamma_x$ to be the curve that passes through every $Z_i$, for $1 \leq i \leq r$, exactly once. This is possible because every pair of zig-zag paths with homology $(0,1)$ is edge-disjoint. This choice of $\gamma_x$ and $\gamma_y$ does not affect $N(\Gamma)$. 

    Following the notation of Proposition \ref{rootsT}, we call the edges of $Z_i$ that follow immediately after the left (resp. right) turn $a$-edges (resp. $b$-edges). In order to describe the marginal polynomial attached to the vertical side of $N(\Gamma)$, we need to choose either all $a$-edges or all $b$-edges of the zig-zag path $Z_r$. In order to traverse the loops of total homology $k/2$ through $\Gamma$, we need to make the same choice on every zig-zag path $Z_i$. For each of them, we can choose either all $a$-edges or all $b$-edges, and by the isoradiality of $\Gamma$, these zig-zag paths are edge-disjoint. See the left part of Figure \ref{fig:Zi_in_Gamma_Gamma_DeltaS}. Therefore, the choice can be made independently for each zig-zag path $Z_i$, $1 \leq i \leq r$. 

    \begin{figure}[h!]
    \centering
    
    \caption{\label{fig:Zi_in_Gamma_Gamma_DeltaS} Left: the portion of the graph $\Gamma$ containing the zig-zag path $Z_i$ with homology $(0,1)$, its $a$-edges are depicted in blue, and $b$-edges are in red. The curve $\gamma_x$ is chosen such that it intersect every zig-zag path $Z_i$ exactly once. Right: the portion of the graph $\Gamma_{\Delta_S}$ obtained by replacing the vertices $v_{i,2}, v_{i,3}, v_{i,6}$ with triangles. The forbidden path from $u_{i,1}$ to $u_{i,4}$ is shown in green.}
    \end{figure}

    We now want to describe the marginal polynomials of $N(\Gamma_{\Delta_S})$. First, notice that by Theorem \ref{equal newton polygons delta}, we have $N(\Gamma_{\Delta_S}) = N(\Gamma)$. If $r = 1$, then the result follows because the marginal polynomial has degree $1$ and is always real-rooted. Therefore, we assume that $r \geq 2$.

    We denote by $V(Z_i) = \{v_{i,1}, \dots, v_{i,k}\}$ the vertex set of $Z_i$ for every $1 \leq i \leq r$. Take $v_{i,j} \in V(\Gamma)$, for $1 \leq j \leq k$. 
    If $v_{i,j} \notin S$, then we denote by $u_{i,j}$ the corresponding vertex in $V(\Gamma_{\Delta_S})$. If $v_{i,j} \in S$, then $u_{i,j}$ is the unique vertex in the triangle $\Delta_{v_{i,j}}$ that is not incident to any edge in $Z_i$ of the initial graph $\Gamma$. See the right part of Figure \ref{fig:Zi_in_Gamma_Gamma_DeltaS}.

    In order to describe the collections of loops in $\Gamma_{\Delta_S}$ whose first coordinate of the homology class is equal to $k/2$, instead of choosing $a$-edges and $b$-edges in $Z_i$, we need to consider the disjoint oriented paths going from $U^-_{i} = \{u_{i,1}, u_{i, 3}, \dots, u_{i, k-1}\}$ to $U^+_{i} = \{u_{i,2}, u_{i, 4}, \dots, u_{i, k}\}$. These paths are chosen such that the rest of the graph can be completed into a collection of loops that is realizable in $\mathcal{L}(\Gamma_{\Delta_S})$.

    It suffices to consider the $k/2$ oriented disjoint paths from $U^-_{i}$ to $U^+_{i}$ and double edges in order to cover all the vertices in $\Delta_{v_{i,j}}$. This covers the part of the graph $\Gamma_{\Delta_S}$ that is obtained from $Z_i$. By Lemma \ref{N(Gamma) < N(Gamma_Delta_v)}, using the original loops of $\Gamma$, we can find appropriate loops in $\Gamma_{\Delta_S}$. We denote by $P_{U^-_{i} \to U^+_{i}}(w)$ the characteristic polynomial describing these oriented paths and double edges.
    
    The only two ways in which these paths can be connected are $u_{i,1} \rightarrow u_{i,2}$, \dots, $u_{i,k-1} \rightarrow u_{i,k}$ or $u_{i,1} \rightarrow u_{i,k}$, \dots, $u_{i,k-1} \rightarrow u_{i,k-2}$. For example, if there is a path from $u_{i,1}$ to $u_{i,4}$, then there is no vertex in $U^-_{i}$ that can be connected to $u_{i,2}$ by a disjoint path. See the right part of Figure \ref{fig:Zi_in_Gamma_Gamma_DeltaS}. With similar reasoning, we can rule out the other options and notice that the only two possibilities are the ones we stated. When all the vertices of $Z_i$ are in $S$, this corresponds to the same argument we had in Lemma \ref{no double edges} c).

    Recall that we chose the curve $\gamma_x$ so that it intersects every $Z_i$ exactly once. By assuming that the added triangles are sufficiently small, $\gamma_x$ intersects only one of the edges in the corresponding portion of the graph $\Gamma_{\Delta_S}$. See Figure \ref{fig:Zi_in_Gamma_Gamma_DeltaS}. This implies that the degree of the polynomial $P_{U^-_{i} \to U^+_{i}}(w)$ is at most one for all $1 \leq i \leq r$. By Theorem \ref{equal newton polygons delta}, it must be equal to one.

    Using similar reasoning as in Lemma \ref{N(Gamma) < N(Gamma_Delta_v)}, we obtain that starting from a collection of loops in $\Gamma$, we can always obtain a collection of loops in $\Gamma_{\Delta_S}$ that uses the two possible collections of paths between $U^-_{i}$ and $U^+_{i}$. Let $P_{ver}(w)$ denote the marginal polynomial describing the vertical side of $N(\Gamma_{\Delta_S})$. We can write it as
    \[P_{ver}(w) = c P_{U^-_{1} \to U^+_{1}}(w) \dots P_{U^-_{r} \to U^+_{r}}(w).\]
    Here, $c$ is the constant that depends on the edge weights of the rest of the graph not contained in any of the paths between $U^-_{i}$ and $U^+_{i}$. This indeed does not depend on the variable $w$, because otherwise, the degree of the polynomial $P_{ver}(w)$ would be greater than $r$, contradicting the fact that $N(\Gamma_{\Delta_S}) = N(\Gamma)$. Since for each $1 \leq i \leq r$, the polynomial $P_{U^-_{i} \to U^+_{i}}(w)$ has a real root, the polynomial $P_{ver}(w)$ is also real-rooted. We can repeat this argument for every side of $N(\Gamma_{\Delta_S})$.
\end{proof}

\end{document}
